\cotefile{145}

\scanpage{20}
\mspage{8}
% archive sheet number: 19
Il revient au même de dire que $\sigma_0,\sigma_1,\sigma_\infty,\rho$ sont des générateurs, avec relations de définition
\[
\overbrace{(\rho^{-1}\sigma_\infty\sigma_0\deleted{\rho})}^{l_\infty}
\overbrace{(\rho^{-1}\sigma_1\sigma_\infty\deleted{\rho})}^{l_1}
\overbrace{(\rho^{-1}\sigma_0\sigma_1\deleted{\rho})}^{l_0}=1.\tag*{(4)}
\]
\[
\sigma_0^2=\sigma_1^2=\sigma_\infty^2=1\tag*{(5)}
\]
\[
\rho^3=1.\tag*{(6)}
\]
% margin note written sideways; slanted
\margin{Il faut ajouter les relations $\rho\sigma_0\rho^{-1}=\sigma_1$ \unsure{etc.} On a deux gén. $\sigma_0,\rho$ en fait avec relations $\sigma_0^2=\rho^3=1$, la \deleted{$[\rho^{-1}\unsure{\sigma_0\rho\sigma_0\rho^{-1}}](\underbrace{\rho\sigma_0\rho^{-1}}_{\sigma_1})(\underbrace{\rho^{-1}\sigma_0\rho}_{\sigma_\infty})\unsure{\rho^{-1}}$} relation (4) étant conséquence des \unsure{précédentes} (\unsure{quand} on y fait $\sigma_1=\rho\sigma_0\rho^{-1}$, $\sigma_\infty=\rho^{-1}\sigma_0\rho$). Noter que $l_0=\deleted{\unsure{(\sigma_0\rho)^2}}\rho^{-1}(\sigma_0\rho)^2\rho=\boxed{\rho^{-1}\sigma_0\rho\sigma_0\rho^{-1}}$, $l_1=(\deleted{\rho}\sigma_0\rho)^2$, $l_\infty=(\rho^{-1}\sigma_0\deleted{\rho})^2$; $\operatorname{int}(\sigma_0)l_1=l_\infty$, $\operatorname{int}(\sigma_0)l_\infty=l_1$, $\operatorname{int}(\sigma_0)(l_0)=l_\infty^{-1}l_1^{-1}$.}
% slanted
\margin{\uline{NB} $\sigma_0(l_0)=l_\infty^{-1}l_1^{-1}$, $\sigma_0(l_1)=l_\infty$, $\sigma_0(l_\infty)=l_1\deleted{\unsure{l_\infty^{-1}}}$ et idem pour $\sigma_1,\sigma_\infty$.}

\uline{NB} \unsure{Si on} y fait \inserted{\uline{de plus}} $\sigma_0,\sigma_1,\sigma_\infty=1$, on trouve comme relation de définition $\rho^3=1$, \deleted{\unsure{alors}} sans plus, ce qui prouve que $\rho$ n'est pas un mot \unsure{p.v.|p.r.} aux $\sigma_i$\unsure{...}

Tout ceci est compatible avec les actions des lacets, \deleted{\unsure{et ces dernières ne seraient pas pour des|et ces dernières ne seraient pas des} \unsure{raisons}, et qu'on interprète $l'_i=l_i^{-1}$, $\sigma'_i=\sigma_{\unsure{\dot\rho|\rho}(i)}$} \deleted{\unsure{puis} $\sigma_i$} \deleted{\unsure{de même que} $\sigma_i'$} si on y remplace les $l_i$ par les $l_i^{-1}$. \unsure{Ici|Déjà} on a, comme $\sigma_i=\sigma'_{\dot\rho^{-1}(i)}$
\[
\left\{
\begin{aligned}
&\operatorname{int}(\sigma_i):\ \text{échange }l_j\text{ et }l_{\deleted{\unsure{k}}h}\quad \inserted{(i,j,h)\text{ distincts}}\text{ et}\\
&\qquad\deleted{\text{(et le dernier, }l_i,\text{ s'en déduit par (1))}}\\
&\operatorname{int}(\rho)l_i=l_{\dot\rho(i)}.
\end{aligned}
\right.\tag*{(7)}
\]

On voudrait maintenant expliciter les $\unsure{\varepsilon_i|\epsilon_i}$ avec $\varepsilon_i^2=l_i$ $(i=0,1,\infty)$, $\varepsilon_i$ étant un générateur du centralisateur (ou normalisateur, c'est \unsure{idem}) des $L_i$ dans $\unsure{\mathscr E^+|\mathcal E^+}$. (Ou le centralisateur dans $\mathscr E$). Les calculs faits plus haut, avec les chgts de notation indiqués, donnent
\[
\left\{\begin{aligned}
\varepsilon_0&=l_0^{\unsure{4|1|-1}}\sigma_1l_{\unsure{1|0|\infty}}^{\unsure{4|1|-1}}\\
\varepsilon_1&=l_\infty^{-1}\sigma_\infty\deleted{l_1^{-1}}l_1\\
\varepsilon_\infty&=l_0^{-1}\sigma_0l_\infty
\end{aligned}\right.\tag*{\unsure{(8)}}
\]

\scanpage{21}
% archive sheet number: 20
\unsure{Donc on a|D'où} les relations
\[
\left\{\begin{aligned}
&\widetilde\sigma_0^2=\widetilde\sigma_1^2=\widetilde\sigma_\infty^2=1\\
&\widetilde\sigma_0\widetilde\sigma_1=\widetilde\sigma_1\widetilde\sigma_\infty=\widetilde\sigma_\infty\widetilde\sigma_0=\deleted{\widetilde\rho}\rho\deleted{(\unsure{=\rho})}\\
&(\deleted{\widetilde\rho}\rho^3=1\text{ est conséquence des précédentes})
\end{aligned}\right.\tag*{(13)}
\]
% slanted
\margin{NB $\rho\widetilde\sigma_i\rho^{-1}=\widetilde\sigma_{\unsure{\dot\rho i|\dot\rho(i)}}$ i.e. $\rho\widetilde\sigma_0\rho^{-1}=\widetilde\sigma_1$, $\rho\widetilde\sigma_1\rho^{-1}=\widetilde\sigma_\infty$, $\rho\widetilde\sigma_\infty\rho^{-1}=\widetilde\sigma_0$.}
et l'image réciproque $\unsure{\widetilde{\mathscr E}|\widetilde{\mathcal E}}$ de $\varphi(\unsure{\widetilde{\mathfrak S}_3|\mathfrak S_3|\widetilde{\mathscr S}_3|\widetilde G_3})\subset\Gamma$ (s-groupe d'indice 2) est produit $\frac12$ direct de $\widetilde{\mathfrak S}_3$ et de $\pi$. Il faut expliciter les opérations des $\widetilde\sigma_a$ sur $\pi$ pour avoir la structure complète du groupe. On trouve que $\widetilde\sigma_a$ est au-dessus de $\varphi(\deleted{\unsure{\widetilde\sigma_a}}\dot\sigma_a)$ i.e. la symétrie \unsure{p.r.|p.v.} au plan passant par $P,\overline P,\deleted{\unsure{a}}\unsure{i}$, qui transforme $l_0$ en $l_0^{-1}$, $\deleted{l}l_\infty$ en $l_1^{-1}$, $l_1$ en $l_\infty^{-1}$ :
% slanted
\margin{\unsure{Ce sont elles obtenues par transport de structure}, via les opérations des \unsure{mêmes} (4,8).}
\[
\left\{\begin{aligned}
\widetilde\sigma_0l_0\widetilde\sigma_0^{-1}&=l_0^{-1}\qquad
\widetilde\sigma_0l_1\widetilde\sigma_0^{-1}=l_\infty^{-1}\qquad
\widetilde\sigma_\infty l_\infty\widetilde\sigma_0^{-1}=l_1^{-1}\\
\widetilde\sigma_1l_1\widetilde\sigma_1^{-1}&=l_1^{-1}\qquad\cdots\\
&\cdots
\end{aligned}\right.\tag*{(14)}
\]
% slanted
\margin{\uline{Critique}. La notation $\widetilde\sigma_0$ suggère à tort que $\widetilde\sigma_0$ est au-dessus de $\sigma_0$, il vaut mieux écrire $\dot\sigma_i'=\dot\sigma_i\tau$ et noter $\widetilde\sigma_i'$ au lieu des $\widetilde\sigma_i$ $(i\in\{0,1,\infty\})$, ou mieux $\sigma_i'$ (sans $\sim$) --- pour réserver $\widetilde\sigma_i$ à des éléments de $\unsure{S\mathscr E|S_\Sigma|S\Sigma}$, cf plus bas\unsure{...}}

Les relations (13), (14), plus la relation fondamentale $l_\infty l_1l_0=1$ (1), font une présentation de $\unsure{\widetilde{\mathscr G}|\widetilde{\mathscr E}|\widetilde G}$, dont on peut si on veut éliminer $\rho$, (en omettant d'écrire la relation qui le définit).

On peut maintenant remonter $\dot\tau$ en trois manières, à savoir
\[
\left\{\begin{aligned}
\tau_0&=\sigma_0\widetilde\sigma_0=\widetilde\sigma_0\sigma_0\\
\tau_1&=\sigma_1\widetilde\sigma_1=\widetilde\sigma_1\sigma_1\\
\tau_\infty&=\sigma_\infty\widetilde\sigma_\infty=\widetilde\sigma_\infty\sigma_\infty
\end{aligned}\right.\tag*{(15)}
\]
\[
\text{donc}\quad
\left\{\begin{aligned}
\tau_0^2&=1\\
\tau_1^2&=1\\
\tau_\infty^2&=1
\end{aligned}\right.\tag*{(16)}
\]

\scanpage{22}
\mspage{9}
% archive sheet number: 21
\[
\left\{\begin{aligned}
&\rho^3=1\\
&\unsure{\varepsilon_\infty|\epsilon_\infty}^{2}\varepsilon_1^2\varepsilon_0^2=1\\
&\deleted{\unsure{l_\infty}}\varepsilon_1=\varepsilon_0^{\unsure{4}}\varepsilon_\infty^4\rho\varepsilon_0^2\deleted{\rho}\varepsilon_\infty^2\\
&\varepsilon_\infty=\varepsilon_1\varepsilon_0^4\rho\varepsilon_{\unsure{\infty|1}}^2\deleted{\rho}\varepsilon_0^2\\
&\varepsilon_0=\varepsilon_\infty\varepsilon_1^4\rho\varepsilon_\infty^2\deleted{\rho}\varepsilon_1^2
\end{aligned}\right.\tag*{(10)}
\]
% slanted
\margin{il faut aussi ajouter (ou \unsure{réintroduire}) les relations $\rho\varepsilon_0\rho^{-1}=\varepsilon_1$ etc donc on a deux gén. $(\varepsilon_0,\rho)$\unsure{!}}

On peut encore se demander si $\rho$ s'exprime en fonction des $\varepsilon_i$, et on trouve comme tantôt que non.

Quant à relever $\dot\tau$ en un élément de $\unsure{\mathscr E|\mathcal E}$, par un chemin de $P$ \unsure{à|:} $\overline P$, on ne voit pas de choix privilégié --- on voit trois choix \unsure{équivalents|également} naturels. Au lieu de les expliciter \unsure{d'emblée|d'abord}, je préfère considérer le sous-groupe d'isotropie de $P$ dans $\Gamma=\Gamma_{0,3}=\unsure{\mathfrak S_3|\mathscr S_3|G_3}\times\mathbb Z/2$, qui est formé des couples $\dot\sigma,\dot\tau^{\varepsilon}$ tels que $\operatorname{sg}(\dot\sigma)=\varepsilon$; \deleted{\unsure{On}} i.e. les éléments
\[
\widetilde{\dot\sigma}=(\dot\sigma,\dot\tau^{\operatorname{sg}(\dot\sigma)})\in\Gamma
\]
qui sont relevés canoniquement en des éléments $(\widetilde{\dot\sigma},\text{ lacet trivial en }P)$ de $\mathscr E$. On trouve un hom.
\[
\left\{\begin{aligned}
\varphi:\ \mathfrak S_3&\longrightarrow\Gamma=\mathfrak S_3\times\mathbb Z/2\\
\dot\sigma&\longmapsto\dot\sigma,\operatorname{sg}(\dot\sigma)
\end{aligned}\right.\tag*{(11)}
\]
et un relèvement canonique de $\varphi$ en
\[
\sigma\longrightarrow\widetilde\sigma:\quad\mathfrak S_3\longrightarrow\mathscr E\tag*{(12)}
\]
On pose
\[
\widetilde\sigma_i=\widetilde{\dot\sigma}_i.
\]
\margin{(on a d'autre part \deleted{$\unsure{\widetilde\rho}$} $\widetilde{\dot\rho}=\rho$\unsure{)}}

\scanpage{23}
\mspage{10}
% archive sheet number: 22
Le sous-groupe $\unsure{\mathscr E_\tau|\mathcal E_\tau}$ de $\mathscr E$ image réciproque de $\mathbb Z/2$ engendré par $\tau$, est le groupe cartographique ``\unsure{triangulé} \unsure{pondéré}'' de Malgoire, ayant comme gén $\tau_0,\tau_1,\tau_\infty$ \deleted{\unsure{relati}} avec
\[
\left\{\begin{aligned}
\deleted{\tau_1\tau_\infty=}{ }l_0&=\tau_\infty\tau_1\\
\deleted{\tau_\infty\tau_0=}{ }l_1&=\tau_0\tau_\infty\\
l_\infty&=\tau_1\tau_0
\end{aligned}\right.\tag*{(17)}
\]
\margin{\uline{NB} \unsure{---}}
$\mathscr E_\tau$ peut être considéré comme le groupe ayant comme \inserted{gén. les $\tau_i$, avec les} seules relations (16) --- on a un hom. signature évident $\tau_i\to-1$ de $\mathscr E_\tau$ dans $\pm1$, dont le noyau est $\pi$ engendré par les $l_i$ décrits par (17)\unsure{....}

\begin{center}---\end{center}
% slanted
\margin{On a les relations}
\margin{$\unsure{\rho}$}
Revenons à une description plus symétrique des autom. extérieurs à lacets de $\mathscr E$. On désignera par $\widehat{\mathscr E}$ le complété profini de $\mathscr E$, extension de $\widehat\Gamma_{0,3}=\Gamma_{0,3}=\unsure{\mathfrak S_3|\mathscr S_3|G_3}\times\mathbb Z/2$ par $\widehat\pi$, et par $\widehat{\widehat{\mathscr E}}$ le groupe des autom. \deleted{extérieurs} à lacets de $\widehat\pi$, extension de $\widehat{\widehat\Gamma}=\operatorname{Aut}\operatorname{ext}_{\unsure{lac|lc}}(\widehat\pi)$ par $\widehat\pi$
\diagram{\(\text{(18)}\quad\)
\begin{tikzcd}[ampersand replacement=\&,column sep=large]
1\arrow[r]\&\widehat\pi\arrow[r]\&\widehat{\widehat{\mathscr E}}\arrow[r]\&\widehat{\widehat\Gamma}\arrow[r]\arrow[d,equal]\&1\\
\&\&\&\operatorname{Aut}\operatorname{ext}_{lac}(\widehat\pi)\&
\end{tikzcd}}

\scanpage{24}
% archive sheet number: 23
qui contient $\unsure{\widehat{\mathscr E}|\widehat{\mathcal E}}$ comme sous-extension. On désigne par $\unsure{\mathbb\Gamma|\Pi}$ le commutant de $\unsure{\mathfrak S_3|\mathscr S_3|G_3}$ dans $\widehat{\widehat\Gamma}$ \inserted{(contenant $\unsure{\dot\tau|\tau}$)}, et par $\unsure{\mathbb E|E}$ l'extension correspondante de $\mathbb\Gamma$ par $\widehat\pi$. Notons déjà que
\[
\left\{\begin{aligned}
\mathbb\Gamma\cap\Gamma_{\unsure{0,3}}&=\{1,\tau\}\\
\mathbb\Gamma\cap\mathfrak S_3&=\{1\}
\end{aligned}\right.
\]
\margin{\deleted{\unsure{NB}}}
% slanted
\margin{\uline{NB} $\widehat{\widehat\Gamma}\supset\mathfrak S_3\times\mathbb\Gamma$; $\mathfrak S_3\times\widehat{\widehat\Gamma}^{!}\supset\widehat{\widehat\Gamma}$. --- j'ignore si cette inclusion est une égalité, i.e. si $\mathbb\Gamma\hookrightarrow\widehat{\widehat\Gamma}^{!}$ est une égalité.}

On veut donner d'abord une description plus symétrique des éléments de $\widehat{\widehat{\mathscr E}}^{!}$, sans faire jouer un rôle particulier à $L_0$ engendré par $l_0$ $(L_0=\widehat k_0(\widehat{\mathbb Z}))$. Un $u\in\widehat{\widehat{\mathscr E}}$ est \inserted{peut être} défini par un multiplicateur
\[
\unsure{p}=\chi(u)\in\widehat{\mathbb Z}^{*}\tag*{(19)}
\]
et des
\[
g_i\in\widehat\pi\qquad i=0,1,\infty\tag*{(20)}
\]
pour
\[
u(l_i)=\operatorname{int}(g_i^{\unsure{-1|1}})l_i^p\tag*{(21)}
\]
les données (19), (20) étant liées par la relation
\[
(\operatorname{int}(g_\infty^{-1})l_\infty^p)(\operatorname{int}g_1^{-1}\cdot l_1^p)(\operatorname{int}(g_0^{-1})l_0^p)=1\tag*{(22)}
\]
et par le fait que les $u(l_i)$ définis par (21) \uline{engendrent} $\widehat\pi$ topologiquement. Pour $u$ fixé, $p$

\scanpage{25}
\mspage{11}
% archive sheet number: 24
est unique, et $g_i$ défini multiplication à gauche par un él. quelconque de $L_i$, soit $l_i^{\alpha_i}$ $(\alpha_i\in\widehat{\mathbb Z})$. On passe de $u$ à $\operatorname{int}(g)\circ u$ en \inserted{gardant $p$, et en} remplaçant les $g_i$ par les $gg_i$.

Si $v$ est donné par $q,(h_i)$, on a
\[
vu(l_i)=v(\operatorname{int}(g_i^{\unsure{-1|1}})l_i^p)=\operatorname{int}(v(g_i^{-1}))\,v(l_i)^p
\]
où
\[
v(l_i)=\operatorname{int}(h_i^{-1})l_i^q\qquad\text{donc}\quad v(l_i)^p=\operatorname{int}(h_i^{-1})l_i^{pq}
\]
donc
\[
vu(l_i)=\operatorname{int}\bigl(\overbrace{v(g_i^{-1})h_i^{-1}}^{\deleted{\unsure{(h_i\cdot v(g_i))^{-1}}}}\bigr)l_i^{pq}.\tag*{(23)}
\]
Si on désigne par $\unsure{S\widehat{\widehat{\mathscr E}}^{!}|S\widehat{\widehat{\mathcal E}}^{!}|S\widetilde{\widehat{\widehat{\mathscr E}}}^{!}|S\widetilde{\widehat{\widehat{\mathcal E}}}^{!}}$ l'extension de $\widehat{\widehat{\mathscr E}}^{!}$ par $\widehat{\mathbb Z}^{\{0,1,\infty\}}$ définie par \deleted{\unsure{les couples}} \unsure{groupe|formés} des \deleted{couples} triples d'un $u\in\widehat{\widehat{\mathscr E}}^{!}$ et d'un système des $g_i$ qui explicitent \inserted{via (21)} \unsure{l'invariance|l'existence} de la structure à lacets ($p$ étant le multiplicateur), on a sur $S\widehat{\widehat{\mathscr E}}^{!}$ trois fonctions à valeurs dans $\deleted{\unsure{\mathscr E}}\widehat\pi$
\[
(u,(g_i))\overset{\gamma_i}{\longmapsto}g_i,
\]
et on a
\[
\begin{aligned}
\gamma_i(\widetilde v\widetilde u)
&=\deleted{\widetilde v(\unsure{\gamma_i(\widetilde u)})\cdot\gamma_i(\widetilde v)}\\
&\quad\deleted{\unsure{\gamma_i(\widetilde v)\cdot v(\gamma_i(\widetilde u))}}\ \inserted{\widetilde v(\gamma_i(\widetilde u))\gamma_i(\widetilde v)},
\end{aligned}\tag*{(24)}
\]
% slanted
\margin{\uline{NB} $S\widehat{\widehat{\mathscr E}}^{!}$ a été construit comme le groupe des systèmes $(p,g_0,g_1,g_\infty)\in\widehat{\mathbb Z}^{*}\times\unsure{\pi|\widehat\pi}^{\{0,1,\infty\}}$ satisfaisant la relation (22).}
i.e. $(\gamma_i)$ est un 1-cocycle sur le groupe \deleted{$\unsure{\widehat\pi}$} \inserted{$S\widehat{\widehat{\mathscr E}}^{!}$} à valeurs dans le groupe \deleted{\unsure{non commutatif}} \inserted{opposé au groupe} $\pi^{\{0,1,\infty\}}$. \deleted{groupe}

Le sous-groupe de $S\widehat{\widehat{\mathscr E}}^{!}$ noyau de $S\widehat{\widehat{\mathscr E}}^{!}\longrightarrow\operatorname{Aut}\operatorname{ext}_{\unsure{lac|lc}}^{!}(\widehat\pi)$ s'identifie au groupe \unsure{pro.} direct $(\widehat\pi\times\prod_{i\in\{0,1,\infty\}}L_i)\simeq\widehat\pi\times\deleted{\unsure{\widehat{\mathbb Z}}}\widehat{\mathbb Z}^{\{0,1,\infty\}}$

\scanpage{26}
% archive sheet number: 25
 en associant à $(g,(\unsure{\alpha^i|\alpha_i})_{i\in\{0,1,\infty\}})$ le système \deleted{\unsure{de}}
\[
p=1,\quad(g_i)=\deleted{\unsure{(l_i^{\alpha_i})}}\deleted{\unsure{(l_i^{\alpha_i}g)}}(gl_i^{\alpha_i})
\]
Il faudrait expliciter les opérations des éléments $\sigma_i$ ou $\unsure{\varepsilon_i|\epsilon_i}$ \deleted{\unsure{sur}} $(\in\unsure{S\widetilde{\widehat{\widehat{\mathscr E}}}^{!}|S\widehat{\widehat{\mathscr E}}^{!}|S\widetilde{\widehat{\widehat{\mathcal E}}}^{!}|S\widehat{\widehat{\mathcal E}}^{!}})$, sur $\deleted{\unsure{\mathscr E}}\widetilde u\inS\widetilde{\widehat{\widehat{\mathscr E}}}^{!}$, \deleted{\unsure{déjà}} par conjugaison, et en particulier $\rho$. On trouve que \deleted{\unsure{(bien sûr)}} les $\gamma_i$ sont ``covariants'' par $\rho$, i.e.
\[
\gamma_i(\rho\widetilde u\rho^{-1})=\gamma_{i+1}(\unsure{\overline u|\widetilde u})
\]
\deleted{\unsure{si} $\widetilde u$ est défini par}
% slanted
\margin{\uline{NB} Il y a lieu d'introduire un groupe $S\unsure{\widetilde{\mathscr E}|\mathscr E}$ extension de $\mathscr E$ par $\unsure{\mathbb Z|\widehat{\mathbb Z}}^{\{0,1,\infty\}}$ (\unsure{évidente d'image extérieure canonique sur|élément d'image extérieure canonique sur} $\mathbb Z^{\{0,1,\infty\}}$) et d'expliciter les relèvements des éléments $\sigma_i,\rho$ en des éléments de $S\widetilde{\mathscr E}$\unsure{...}}

\scanpage{27}
\mspage{12}
% archive sheet number: 26
L'extension $\unsure{\mathscr E|\mathcal E}=\mathscr E_{0,3}$ de
\[
\Gamma=\Gamma_{0,3}\simeq\underbrace{\unsure{\mathfrak S_{\{0,1,\infty\}}|\mathscr S_{\{0,1,\infty\}}|G_{\{0,1,\infty\}}}}_{=\mathfrak S_3}\times\mathbb Z/2\mathbb Z
\]
\margin{engendré par $\tau$: $z\mapsto\overline z$}
par
\[
\pi=\pi_{0,3}=\{l_0,l_1,l_\infty\mid l_\infty l_1l_0=1\}=\pi_1(\underbrace{\mathbb P^1_{\mathbb C}\setminus\{0,1,\infty\}}_{\unsure{U=X_{0,3}|U=U_{0,3}}},-\overline j\}.
\]
Les classes de conjugaison \deleted{\unsure{au dessus des|au dessus de}} de scindages partiels maximaux correspondent aux points fixes de $\dot\rho,\dot\sigma_0,\dot\sigma_1,\dot\sigma_\infty$ opérant sur $X_{0,3}$, et aux comp. connexes des ens. d'invariants \inserted{(éventuelles)} de $\deleted{\dot\rho}\dot\tau,\dot\sigma_0\tau,\dot\sigma_1\tau,\dot\sigma_\infty\tau$. Les pts fixes de $\rho$ sont $P$ et $\overline P$ (dont le stabilisateur est d'ordre 6, \inserted{savoir $\varphi(\mathfrak S_3)$}) formé de $1,\dot\rho,\dot\rho^2,\dot\sigma_1\tau,\dot\sigma_2\tau,\dot\sigma_3\tau$. On a étudié le relèvement correspondant \inserted{à $P$} (de $\varphi(\mathfrak S_3)$), \unsure{on trouve} des relèvements ($\pi$-conjugués \inserted{entre eux}) relatifs à $\overline P$ s'obtiennent en \deleted{\unsure{conj}} prenant un élément de $\mathscr E$ \deleted{\unsure{et}} qui soit au dessus d'un élément de $\Gamma$ envoyant $P$ en $\overline P$ (savoir $\dot\sigma_0,\dot\sigma_1,\dot\sigma_\infty,\tau,\tau\rho,\tau\rho^2$) et en conjuguant par cet élément on peut donc prendre \deleted{$\unsure{\sigma_0}$} un des éléments $\sigma_0,\sigma_1,\sigma_\infty$ (ou $\unsure{\varepsilon_0|\epsilon_0},\varepsilon_1,\varepsilon_\infty$) ou $\tau_0,\tau_1,\tau_\infty$ ou les produits des précédents par \deleted{\unsure{un}} $\rho$ ou $\rho^2$\unsure{...}.

Le \inserted{seul} pt fixe de $\sigma_0$ est 2, dont le stabilisateur est le groupe $\simeq\mathbb Z/2\times\mathbb Z/2$ engendré par $\dot\sigma_0$ et $\tau$. On a un relèvement de ce sous-groupe

\scanpage{28}
% archive sheet number: 27
dans $\unsure{\mathscr E|\mathcal E}$, qu'on peut choisir canoniquement en prenant le quart de grand cercle qui joint $P$ au milieu 2 du \deleted{segment} l'arc $(1,\infty)$ de grand cercle \unsure{opposée|opposé} à 0; ce qui définit un isom du rev universel de $\unsure{U|X}$ en $P$, \unsure{avec} celui de $U$ en $\unsure{\overline P|2}$ sur lequel le groupe précédent opère de façon canonique, d'où aussi par transport de structure sur $\widetilde U(\deleted{\unsure{\overline P}}P)$, donc il s'envoie dans $\mathscr E$.

On trouve que
\[
\dot\sigma_0\mapsto\sigma_0,\quad\dot\tau\to\tau_0
\]
est le remontage correspondant. De $=$ pour les pts fixes $-1$ de $\dot\sigma_1$, ou $\frac12$ de $\dot\sigma_\infty$, on trouve des scindages partiels sur les ss-groupes des \deleted{\unsure{groupes}} \unsure{stabs} $\Gamma_{-1}\subset\Gamma$ et $\Gamma_2\subset\Gamma$. Il est commode de poser
\[
\deleted{\unsure{2}}\quad 2=0',\qquad -1=1',\quad 1/2=\infty',
\]
d'où des groupes d'isotropie $\Gamma_{0'}=\{\dot\sigma_0,\tau\}$, $\Gamma_{1'}=\{\dot\sigma_1,\tau\}$, et $\Gamma_{\infty'}=\{\dot\sigma_\infty,\tau\}$, qui se remontent resp en $\sigma_0$ et $\tau_0$, en $\sigma_1$ et $\tau_1$, en $\sigma_\infty$ et $\tau_\infty$.

Les composantes connexes de $U^\tau$ sont les trois arcs de grand cercle entre $0,1,\infty$, qui contiennent chacun un \deleted{pt} et un seul des pts $0',1',\infty'$, et les classes de conjugaison

\scanpage{29}
\mspage{13}
% archive sheet number: 28
correspondantes de relèvements \inserted{de $\{1,\dot\tau\}$} sont resp. les relèvements en $\tau_0$, en $\tau_1$ et en $\tau_\infty$ --- qui ne donnent rien de nouveau. Reste à considérer $\tau\dot\sigma_1$ (et ses conjugués \inserted{par $\rho$} $\tau\dot\sigma_\infty$, $\unsure{\tau\dot\sigma_0|\tau_\infty\dot\sigma_0}$), on a $\dot\sigma_1(2)=1/2$, donc
\[
(\tau\dot\sigma_1)(2)=\frac12
\]
et
\[
U^{\tau\dot\sigma_1}=\{z\in U\mid z\overline z=1\}=\mathbb U-\{1\},
\]
où $\mathbb U=\{z\in\mathbb C\mid |z|=1\}$. Cet ensemble contient $P$ et $\overline P$, et ne donne donc pas de \unsure{nouveaux|nouveau} \unsure{relèvements|relèvement} \unsure{p.v.|p.r.} à celui déjà donné, avec $\widetilde\sigma_1$ relevant $\dot\sigma_1\tau$.

Il faudrait quand-même donner aussi une description de $\unsure{\mathscr E|\mathcal E}$ \unsure{tt|tout} entier par gén. et relations, en partant d'un syst. de gén. de $\Gamma=\unsure{\mathfrak S_3|\mathscr S_3|G_3}\times\mathbb Z/2$, qu'on va relever de façon judicieuse. \deleted{On prendra les \inserted{quatre} générateurs}
\[
\deleted{\unsure{\rho'|\rho}=\rho\dot\tau,\quad
\underset{\dot\sigma_0\tau}{\overset{\widetilde{\dot\sigma}_0}{\text{''}}},\quad
\underset{\dot\sigma_1\tau}{\overset{\widetilde{\dot\sigma}_1}{\text{''}}},\quad
\underset{\dot\sigma_\infty\tau}{\overset{\widetilde{\dot\sigma}_\infty}{\text{''}}}}
\]
\deleted{\inserted{(donc $\rho'^{\,6}=1$, $\rho'^{\,3}=\dot\tau$, $\rho'^{\,2}=\rho^2=\rho^{-1}$)}}
\deleted{avec les relations}
\[
\deleted{\left\{\begin{aligned}
\widetilde{\dot\sigma}_1^2&=\widetilde{\dot\sigma}_2^2=\widetilde{\dot\sigma}_3^2=1\\
\widetilde{\dot\sigma}_0\widetilde{\dot\sigma}_1&=\widetilde{\dot\sigma}_1\widetilde{\dot\sigma}_\infty=\widetilde{\dot\sigma}_\infty\widetilde{\dot\sigma}_0=\unsure{\rho'^{\,2}|\rho^2}\\
\rho'^{\,6}&=1\\
\rho'^{\,3}&\text{ commute aux }\dot\sigma_i
\end{aligned}\right.}
\]

\scanpage{30}
% archive sheet number: 29
\deleted{qui décrivent $\unsure{\mathfrak S_3|\mathscr S_3|G_3}\times\mathbb Z/2$, on va relever ces générateurs en des éléments de $\unsure{\mathscr E|\mathcal E}$, de façon que $l_\infty,l_1,l_0$ s'expriment comme des mots en ces éléments, et il \unsure{suffira|suffit} alors d'ajouter cette relation $l_\infty l_1l_0$ aux relations correspondant aux précédentes\unsure{....}}

$\dot\sigma_0,\dot\sigma_1,\dot\sigma_\infty$, $\unsure{\dot\rho|\rho},\tau$, avec les relations
\[
\left\{\begin{aligned}
\dot\sigma_0^2&=\dot\sigma_1^2=\dot\sigma_\infty^2=1\\
\dot\sigma_0\dot\sigma_1&=\dot\sigma_1\dot\sigma_\infty=\dot\sigma_\infty\dot\sigma_0=\dot\rho\\
\dot\rho^3&=1\\
\tau^2&=1\\
\tau&\text{ commute à }\dot\sigma_0,\dot\sigma_1,\dot\sigma_\infty\quad(\text{donc aussi à }\dot\rho)
\end{aligned}\right.
\]
on les relève en $\deleted{\tau_0,\tau_1,\tau_\infty}\sigma_0,\sigma_1,\sigma_\infty,\rho,\underbrace{\tau_0,\tau_1,\tau_\infty}_{\text{sur }\tau}$ avec les relations
\[
\left\{\begin{aligned}
\sigma_0^2&=\sigma_1^2=\sigma_\infty^2=\tau_0^2=[\tau_1^2=\tau_\infty^2]=1\\
\rho^3&=1\\
\rho\tau_0\rho^{-1}&=\tau_1,\quad\rho\tau_1\rho^{-1}=\tau_\infty,\quad\rho\tau_\infty\rho^{-1}=\tau_0\\
\tau_0\sigma_0&=\sigma_0\tau_0,\quad\tau_1\sigma_1=\sigma_1\tau_1,\quad\tau_\infty\sigma_\infty=\sigma_\infty\tau_\infty\\
\sigma_0\sigma_1&=\tau_\infty\tau_1\rho,\quad\sigma_1\sigma_\infty=\tau_0\tau_\infty\rho,\quad\sigma_\infty\sigma_0=\tau_1\tau_0\rho
\end{aligned}\right.
\]
% slanted
\margin{\deleted{\unsure{NB}} \unsure{Il faut voir si|Il faut vérifier si} les relations $\left\{\begin{aligned}\rho\sigma_0\rho^{-1}&=\sigma_1\\\rho\sigma_1\rho^{-1}&=\sigma_\infty\\\rho\sigma_\infty\rho^{-1}&=\sigma_0\end{aligned}\right.$ sont conséquences des relations ci-contre.}
\margin{Ce sont les éléments $\widetilde\sigma_0,\widetilde\sigma_1,\widetilde\sigma_\infty$ $\longrightarrow$}
\margin{$\bigl[$ NB \unsure{On trouve} $\begin{aligned}l_0&=\tau_\infty\tau_1,\quad\deleted{l_1=\tau_0\tau_\infty,}\\l_1&=\tau_0\tau_\infty\deleted{\unsure{l_\infty=}}\\l_\infty&=\tau_1\tau_0\end{aligned}$ et $l_\infty l_1l_0=1$ est automatique $\bigr]$}
\margin{\deleted{\unsure{NB}} Ceci implique que $\tau_0,\tau_1,\tau_\infty$ \unsure{conjuguent|engendrent} \unsure{le|un} ss-groupe engendré par $l_0,l_1,l_\infty$ $\uparrow$}

\deleted{\uline{Question} \uline{\deleted{Peut-on}\inserted{On peut}} exprimer $\rho$ en fonction de $\sigma_0,\sigma_1,\sigma_\infty,\tau_0,\tau_1,\tau_\infty$ ??, \deleted{d'après les} \inserted{l'une des} derniers groupes de relations, et substituer dans les autres.}

\scanpage{31}
\mspage{14}
% archive sheet number: 30
On peut prendre aussi comme seuls gén. $\rho,\sigma_0,\tau_0$ (dont se déduisent $\sigma_1,\tau_1$, \inserted{$\sigma_\infty,\tau_\infty$} comme $\rho\sigma_0\rho^{-1}$, $\rho\tau_0\rho^{-1}$, $\rho^2\sigma_0\rho^{-2}$, $\rho^2\tau_0\rho^{-2}$) et les relations
\[
\left\{\begin{aligned}
\unsure{\mathrm a)|1)}\quad&\sigma_0^2=\tau_0^2=1\quad\tau_0\sigma_0=\sigma_0\tau_0\quad(\text{i.e. }(\overbrace{\sigma_0\tau_0}^{\widetilde\sigma_0})^2=1)\\
\unsure{\mathrm b)|2)}\quad&\rho^3=1\\
&\sigma_0\underbrace{\rho\sigma_0^{-1}\rho^{-1}}_{\sigma_1}=\underbrace{\rho^2\tau_0\rho^{-2}}_{\tau_\infty}\underbrace{\rho\tau_0\rho^{-1}}_{\tau_1}\rho\quad\text{i.e. }\sigma_0\rho\sigma_0^{-1}\rho^{-1}=\rho^{-1}\tau_0\rho\tau_0\\
&\inserted{\text{ou enfin}}\quad\rho\sigma_0\rho^{-1}\sigma_0=\tau_0\rho^{-1}\tau_0\rho\quad\text{soit}\quad\deleted{\operatorname{int}(\rho)\sigma_0}\ [\rho,\sigma_0]=[\tau_0,\rho^{-1}]\\
&\text{soit}\quad[\rho,\sigma_0][\rho^{-1},\tau_0]=1
\end{aligned}\right.
\]
\margin{\deleted{i.e.} \unsure{ou encore}}

C'est donc un quotient du produit libre $(\mathbb Z/2\mathbb Z\times\mathbb Z/2\mathbb Z)*\mathbb Z/3\mathbb Z$, par le sous-groupe invariant décrit par la relation $[\rho,\sigma_0][\rho^{-1},\tau_0]=1$.

\margin{\uline{NB}}
On trouve une symétrie dans le groupe, qui échange $\sigma_0$ et $\tau_0$, et échange $\rho$ et $\rho^{-1}$. [donc qui échange aussi $\sigma_1$ et $\tau_\infty$, $\sigma_\infty$ et $\tau_1$] --- mais cette symétrie n'est pas compatible avec la structure d'extension, comme on voit \unsure{tout} de suite (en en passant au quotient $\unsure{\mathfrak S_3|\mathscr S_3|G_3}\times\mathbb Z/2$, elle devrait conserver le centre!).

\scanpage{32}
% archive sheet number: 31
Il faudrait encore expliciter l'extension $S\Gamma=S\Gamma_{0,3}$ de $\Gamma$ par $\mathbb Z^{\{0,1,\infty\}}$, ce qui nous donnera (par produits fibrés) une description de $S\unsure{\mathscr E|\mathcal E}=S\Gamma\times_\Gamma\mathscr E$.

\scanpage{33}
\mspage{15}
% archive sheet number: 32
\textcircled{2} \uline{Le groupe} $S\unsure{\widehat{\widehat{\mathscr E}}|\widehat{\widehat{\mathcal E}}}$.

Par le calcul, il s'explicite comme l'ens des systèmes
\diagram{$(1)\qquad$
\begin{tikzcd}[ampersand replacement=\&,row sep=small,column sep=small]
(p,\arrow[d,phantom,"\in"{rotate=-90}]\&\sigma,\arrow[d,phantom,"\in"{rotate=-90}]\&g_0,\arrow[d,phantom,"\in"{rotate=-90}]\&g_1,\arrow[d,phantom,"\in"{rotate=-90}]\&g_\infty)\arrow[d,phantom,"\in"{rotate=-90}]\&=\widetilde u\\
\widehat{\mathbb Z}^{*}\&\unsure{\mathfrak S_3|\mathscr S_3|G_3}\&\widehat\pi\&\widehat\pi\&\widehat\pi\&
\end{tikzcd}}
satisfaisant la condition
\[
\operatorname{int}(g_{\unsure{\infty|0}})(l_{\sigma(\unsure{\infty|0})}^{p})\cdot\operatorname{int}(g_1)(l_{\sigma(1)}^{p})\cdot\operatorname{int}(g_{\unsure{0|\infty}})(l_{\sigma(\unsure{0|\infty})}^{p})=1\tag*{(2)}
\]
et tels que les trois facteurs de ce produit soient un syst. de générateurs de $\widehat\pi$ [avec la relation génératrice (2)]. Le morphisme $u$ défini par ce \deleted{\unsure{morphi}} système est
\[
\left\{\begin{aligned}
u(l_0)&\unsure{\Longrightarrow|\longmapsto|=|\longrightarrow}\operatorname{int}(g_0)l_{\sigma(0)}^p\\
u(l_1)&=\operatorname{int}(g_1)l_{\sigma(1)}^p\\
u(l_\infty)&=\operatorname{int}(g_\infty)l_{\sigma(\infty)}^p
\end{aligned}\right.\tag*{(3)}
\]
\inserted{ce qui explicite l'homom}
\diagram{$(4)\qquad$
\begin{tikzcd}[ampersand replacement=\&]
S\widehat{\widehat{\mathscr E}}\arrow[r,"p"]\&\widehat{\widehat{\mathscr E}}=\operatorname{Aut}_{\mathrm{lac}}(\widehat\pi)
\end{tikzcd}}
L'homomorphisme canonique
\diagram{$(5)\qquad i:\quad$
\begin{tikzcd}[ampersand replacement=\&]
\widehat\pi\times\widehat{\mathbb Z}^{\{0,1,\infty\}}\arrow[r]\&S\widehat{\widehat{\mathscr E}}
\end{tikzcd}}
s'explicite par
\[
i(g;\alpha_0,\alpha_1,\alpha_\infty)=(1,1,g l_0^{\alpha_0},g l_{\unsure{1|0}}^{\alpha_1},g l_\infty^{\alpha_\infty})\qquad\text{donc}\tag*{(6)}
\]
\[
\deleted{\text{donc}}\quad p\,i(g;\alpha_0,\alpha_1,\alpha_\infty)=\operatorname{int}g.\tag*{(7)}
\]
% slanted
\margin{On retrouve que \unsure{$(p)$|$(4)$} est surjectif, et a comme noyau $i(\widehat\pi)$.}
La multiplication dans $\widehat{\widehat{\mathscr E}}$ est donnée par
\[
\underbrace{(p,\sigma,g_0,g_1,g_\infty)}_{\widetilde u}\underbrace{(p',\sigma',g'_0,g'_1,g'_\infty)}_{\widetilde u'}=(pp',\sigma\sigma',u(g'_0)g_{\sigma'(0)},u(g'_1)g_{\sigma'(1)},u(g'_\infty)g_{\sigma'(\infty)}\tag*{(8)}
\]
% The right edge of (8) does not show a secure final tuple parenthesis.
Cette formule n'est pas, du point de vue du calcul

\scanpage{34}
% archive sheet number: 33
entièrement explicite et satisfaisante, puisque au deuxième membre interviennent les quantités $u(g'_0),u(g'_1),u(g'_\infty)$, qui ne sont donc pas exprimées directement en termes des \inserted{seules} données $(p,\sigma,g_0,g_1,g_\infty);(p',\sigma',g'_0,g'_1,g'_\infty)$; \deleted{sauf} \inserted{mais ils s'explicitent \unsure{ainsi|aussi}} quand $g'_0,g'_1,g'_\infty$ sont eux-mêmes donnés explicitement comme des mots (\unsure{ou|en} limites de tels) en les $l_0,l_1,l_\infty$.

Cependant quand $\widetilde u$ est de la forme $i(g,\alpha_0,\alpha_1,\alpha_\infty)$, on trouve (en laissant tomber les $'$ du $2^e$ facteur)
\[
i(g,\alpha_0,\alpha_1,\alpha_\infty)(p,\sigma,g_0,g_1,g_\infty)=(p,\sigma,\deleted{\unsure{\operatorname{int}}}g g_0l_{\sigma0}^{\alpha_{\sigma(0)}}\deleted{\unsure{g^{-1}|g_0}},g g_1l_{\sigma1}^{\alpha_{\sigma(1)}},g g_\infty l_{\sigma\infty}^{\alpha_{\sigma\infty}}\unsure{)}\tag*{(9)}
\]
qui se simplifie pour $\sigma=1$ en
\[
i(g,\alpha_0,\alpha_1,\alpha_\infty)(p,1,g_0,g_1,g_\infty)=(p,\sigma,g g_0l_0^{\alpha_0},g g_1l_1^{\alpha_1},g g_\infty l_\infty^{\alpha_\infty})\tag*{(9')}
\]
Par ailleurs, l'opération de $\unsure{\widehat{\widehat{\mathscr E}}|\widehat{\widehat{\mathcal E}}}$ sur $i(\pi\times\mathbb Z^{\{0,1,\infty\}})$ par autom int. s'explicite aisément, et ne dépend que de $u$ \deleted{\unsure{et}}.
\[
\widetilde u\,i(g;\alpha_0,\alpha_1,\alpha_\infty)\deleted{u}\widetilde u^{-1}=i(u(g);p\alpha_{\sigma^{-1}0},p\alpha_{\sigma^{-1}1},p\alpha_{\sigma^{-1}\infty})\tag*{(10)}
\]
où $\widetilde u=(p,\sigma,g_0,g_1,g_\infty)$

--- mais sur cette formule, on peut faire les mêmes commentaires que plus haut, sur le manque d'explicitation du \unsure{$2^{\text{ième}}$ mbr|$2^{\text{ème}}$ mbr|deuxième mbr} de cette formule, pour $g$ non exprimé comme un mot \unsure{p.v.|p.r.} aux $l_0,l_1,l_\infty$ (ou limites de tels mots)
